\section{Аппроксимация функциональных зависимостей}
\label{sec:function-approximation}

\subsection{Постановка задачи}

Аппроксимацией называется замена функции $f$ более простой функцией $\varphi$ из заранее выбранного класса $V$:
\[
f(x)\approx \varphi(x), \qquad \varphi\in V.
\]

Исходная зависимость может быть задана аналитически, таблицей значений $(x_i,y_i)$ или экспериментальными данными. Класс аппроксимирующих функций
выбирают с учётом требуемой точности, гладкости и вычислительной сложности.

Следует различать:
\begin{itemize}
	\item \textbf{интерполяцию}, при которой $\varphi(x_i)=y_i$ во всех заданных узлах;
	\item \textbf{среднеквадратичную аппроксимацию}, при которой минимизируется средняя квадратичная ошибка;
	\item \textbf{равномерную аппроксимацию}, при которой минимизируется максимальное отклонение функции от приближения.
\end{itemize}

При наличии ошибок измерений обычно не требуется точного прохождения аппроксимирующей функции через все экспериментальные точки.

\subsection{Нормы и критерии качества}

Качество приближения характеризуют величиной ошибки
\[
e(x)=f(x)-\varphi(x).
\]

На отрезке $[a,b]$ часто используют равномерную норму
\[
\|e\|_{\infty}
=
\max_{x\in[a,b]}|e(x)|
\]
и среднеквадратичную норму
\[
\|e\|_2
=
\left(
\int_a^b |e(x)|^2\,dx
\right)^{1/2}.
\]

Для таблично заданной функции соответствующий дискретный критерий имеет вид
\[
\|e\|_{2,h}
=
\left(
\sum_{i=1}^{N}
w_i
\bigl(y_i-\varphi(x_i)\bigr)^2
\right)^{1/2},
\]
где $w_i>0$ --- веса отдельных наблюдений.

При $w_i=1$ все точки считаются равноценными. Если точность измерений различна, веса могут выбираться с учётом погрешности исходных данных.

\subsection{Основные классы аппроксимирующих функций}

\begin{enumerate}
	\item \textbf{Алгебраические многочлены}
	\[
	\varphi(x)
	=
	c_0+c_1x+\ldots+c_mx^m.
	\]
	
	Они просты в использовании, однако применение многочленов высокой степени может приводить к плохой обусловленности и сильным колебаниям приближения.
	
	\item \textbf{Разложения по ортогональным функциям}
	\[
	\varphi_m(x)
	=
	\sum_{j=0}^{m}c_j\phi_j(x),
	\]
	где функции $\phi_j$ ортогональны относительно выбранного скалярного произведения. Примерами являются многочлены Лежандра и Чебышёва, а также тригонометрические функции.
	
	Степенной базис $1,x,x^2,\ldots$ при больших степенях часто плохо обусловлен. Ортогональные базисы обычно численно предпочтительнее.
	
	\item \textbf{Сплайны} --- кусочно-полиномиальные функции, согласованные в точках стыка с заданной степенью гладкости.
	
	Наиболее распространены кубические сплайны. Их важные преимущества перед одним многочленом высокой степени --- малая степень на каждом участке и возможность задавать гладкость в узлах. Следует различать \emph{локальность базиса} и \emph{локальность интерполянта}: базисные B-сплайны имеют компактные носители, но коэффициенты классического интерполирующего кубического сплайна обычно определяются из общей трёхдиагональной системы, поэтому изменение одного значения данных в общем случае может повлиять на весь интерполянт.
	
	\item \textbf{Рациональные функции}
	\[
	\varphi(x)=\frac{P_m(x)}{Q_n(x)}.
	\]
	
	Они могут давать хорошее приближение функциям с резкими изменениями, асимптотическим поведением или особенностями.
	
	\item \textbf{Экспоненциальные и степенные функции}
	\[
	y=ae^{bx},
	\qquad
	y=ax^b.
	\]
	
	Такие зависимости часто возникают непосредственно из физической постановки задачи. В некоторых случаях их параметры можно определять после соответствующего преобразования переменных.
\end{enumerate}

\subsection{Аппроксимация линейной комбинацией базисных функций}

Пусть приближение ищется в виде
\[
\varphi(x)
=
\sum_{j=0}^{m} c_j\phi_j(x),
\]
где $\phi_j(x)$ --- заданные базисные функции, а коэффициенты $c_j$ неизвестны.

Для табличных данных $(x_i,y_i)$ получаем систему
\[
\sum_{j=0}^{m}c_j\phi_j(x_i)\approx y_i,
\qquad
i=1,\ldots,N.
\]

Введём матрицу
\[
A_{ij}=\phi_j(x_i).
\]
Тогда задача записывается в виде
\[
Ac\approx y.
\]

Если $N>m+1$, система обычно переопределена, поэтому точного решения
может не существовать. В методе наименьших квадратов выбирают коэффициенты,
минимизирующие сумму квадратов отклонений:
\[
\min_c \|Ac-y\|_2^2.
\]

Условие минимума приводит к нормальной системе
\[
A^{T}Ac=A^{T}y.
\]

Геометрически вектор $Ac$ является ортогональной проекцией вектора данных
$y$ на линейную оболочку столбцов матрицы $A$.

\subsection{Наилучшее приближение и принцип ортогональности}

Пусть $V_m$ --- конечномерное подпространство нормированного пространства.
Элемент $\varphi_*\in V_m$ называется элементом наилучшего приближения к $f$, если
\[
    \|f-\varphi_*\|
    =\inf_{\varphi\in V_m}\|f-\varphi\|.
\]
В гильбертовом пространстве для замкнутого подпространства $V_m$ такое
приближение существует и единственно; оно является ортогональной проекцией:
\[
    (f-\varphi_*,v)=0
    \qquad\forall v\in V_m.
\]
Именно это условие в дискретном МНК приводит к нормальным уравнениям
$A^T(Ac-y)=0$.

Для равномерного приближения многочленами центральным результатом является
\textbf{теорема Чебышёва об альтернансе}. Многочлен $P_n$ степени не выше $n$
является многочленом наилучшего равномерного приближения непрерывной функции
$f$ на $[a,b]$ тогда и только тогда, когда существуют по крайней мере $n+2$
точки
\[
    a\le x_0<x_1<\cdots<x_{n+1}\le b,
\]
в которых ошибка достигает своей максимальной по модулю величины с
чередующимися знаками:
\[
    f(x_i)-P_n(x_i)
    =\sigma(-1)^i\|f-P_n\|_\infty,
    \qquad \sigma\in\{1,-1\}.
\]
Такой многочлен единствен.

\subsection{Теорема Вейерштрасса}

Теорема Вейерштрасса утверждает, что для любой функции $f$, непрерывной на отрезке $[a,b]$, и любого $\varepsilon>0$ существует многочлен $P(x)$ такой, что
\[
\max_{x\in[a,b]}|f(x)-P(x)|<\varepsilon.
\]

Таким образом, любую непрерывную на отрезке функцию можно с любой заданной точностью равномерно приблизить многочленом.

Однако теорема носит теоретический характер и не утверждает, что увеличение степени многочлена всегда является хорошим численным способом построения приближения. На практике многочлены высокой степени могут быть плохо обусловлены, поэтому часто используют ортогональные многочлены или сплайны.

\subsection{Выбор вида и сложности аппроксимации}

При выборе аппроксимирующей функции необходимо учитывать:
\begin{itemize}
	\item величину ошибки приближения;
	\item гладкость исходной зависимости;
	\item устойчивость решения к погрешностям исходных данных;
	\item обусловленность возникающей системы уравнений;
	\item сложность вычисления приближающей функции;
	\item наличие априорной информации о поведении исследуемой зависимости.
\end{itemize}

Неоправданное увеличение числа параметров не всегда улучшает результат: аппроксимация может становиться чувствительной к ошибкам исходных данных. Поэтому предпочтительно выбирать наиболее простой класс функций, обеспечивающий требуемую точность.

\subsection{Ортогональная проекция и обобщённые ряды Фурье}

Пусть $H$ --- гильбертово пространство, а $V_m=\operatorname{span}\{\phi_0,\ldots,\phi_m\}$ ---
конечномерное подпространство. Наилучшее среднеквадратичное приближение $P_{V_m}f$
характеризуется условием ортогональности
\[
    (f-P_{V_m}f,v)=0\qquad \forall v\in V_m.
\]
Если базис $\{\phi_k\}$ ортогонален, то коэффициенты вычисляются независимо:
\[
    P_{V_m}f=\sum_{k=0}^{m}c_k\phi_k,
    \qquad
    c_k=\frac{(f,\phi_k)}{(\phi_k,\phi_k)}.
\]
Для ортонормированного базиса $c_k=(f,\phi_k)$. При этом выполняется теорема Пифагора
\[
    \|f\|^2=\|P_{V_m}f\|^2+\|f-P_{V_m}f\|^2,
\]
поэтому увеличение подпространства не увеличивает ошибку наилучшего приближения.
Для полной ортонормированной системы предельный переход приводит к разложению Фурье
в соответствующем гильбертовом пространстве.

\subsection{Смысл теоремы Чебышёва и алгоритм Ремеза}

Теорема об альтернансе не только даёт критерий оптимальности, но и объясняет структуру
наилучшего равномерного приближения: ошибка должна ``уравновешиваться'' на отрезке,
достигая одинаковой максимальной величины с чередующимся знаком. Если таких точек
меньше $n+2$, можно построить малое возмущение многочлена, уменьшающее максимальную
ошибку, что противоречит оптимальности.

На этом основан \textbf{алгоритм Ремеза}. Выбираются $n+2$ точки и решается система
\[
    f(x_i)-P_n(x_i)=(-1)^i\sigma E,
    \qquad i=0,\ldots,n+1,
\]
относительно коэффициентов $P_n$ и неизвестной величины $E$. Затем по фактическим
экстремумам ошибки набор точек обновляется. При сходимости получается альтернирующий
набор и, следовательно, многочлен наилучшего равномерного приближения.

\subsection{Устойчивость аппроксимации и выбор базиса}

Две аппроксимации одинаковой теоретической точности могут сильно различаться численно.
Например, матрица Вандермонда в степенном базисе становится плохо обусловленной при
росте степени. Поэтому на практике используют масштабирование переменной,
ортогональные многочлены, QR-разложения, B-сплайны и другие устойчивые базисы.

Следует различать \emph{ошибку аппроксимации} (насколько выбранный класс функций способен
приблизить исходную зависимость) и \emph{вычислительную ошибку} определения коэффициентов.
Увеличение размерности пространства обычно уменьшает первую, но может увеличить вторую.
Именно этот компромисс определяет разумную сложность приближающей модели.

\subsection{Что важно сказать на экзамене}

Аппроксимация состоит в замене исходной зависимости функцией из некоторого более простого класса. В отличие от интерполяции, точное совпадение со всеми заданными значениями обычно не требуется.

Необходимо задать класс аппроксимирующих функций и критерий ошибки. Для приближения линейной комбинацией базисных функций при квадратичном критерии возникает задача метода наименьших квадратов.

Основными средствами аппроксимации являются многочлены, ортогональные разложения,
сплайны и рациональные функции. Для среднеквадратичного критерия ключевым является
принцип ортогональной проекции, для наилучшего равномерного приближения
многочленами --- теорема Чебышёва об альтернансе. Возможность равномерно
приближать любую непрерывную функцию многочленами обеспечивает теорема
Вейерштрасса.

\newpage